%----------
%	APPENDIX X: ZERO SHOT - PROMPTS FOR CONTENIDO NEGATIVO
%----------	

\lstset{
  basicstyle=\ttfamily\scriptsize,
  breaklines=true,
  breakatwhitespace=true,
  columns=fullflexible,
  frame=single,
  inputencoding=utf8,
  extendedchars=true,
  literate=%
    {á}{{\'a}}1 {é}{{\'e}}1 {í}{{\'i}}1 {ó}{{\'o}}1 {ú}{{\'u}}1
    {Á}{{\'A}}1 {É}{{\'E}}1 {Í}{{\'I}}1 {Ó}{{\'O}}1 {Ú}{{\'U}}1
    {à}{{\`a}}1 {è}{{\`e}}1 {ì}{{\`i}}1 {ò}{{\`o}}1 {ù}{{\`u}}1
    {À}{{\`A}}1 {È}{{\`E}}1 {Ì}{{\`I}}1 {Ò}{{\`O}}1 {Ù}{{\`U}}1
    {ä}{{\"a}}1 {ë}{{\"e}}1 {ï}{{\"i}}1 {ö}{{\"o}}1 {ü}{{\"u}}1
    {Ä}{{\"A}}1 {Ë}{{\"E}}1 {Ï}{{\"I}}1 {Ö}{{\"O}}1 {Ü}{{\"U}}1
    {â}{{\^a}}1 {ê}{{\^e}}1 {î}{{\^i}}1 {ô}{{\^o}}1 {û}{{\^u}}1
    {Â}{{\^A}}1 {Ê}{{\^E}}1 {Î}{{\^I}}1 {Ô}{{\^O}}1 {Û}{{\^U}}1
    {ã}{{\~a}}1 {õ}{{\~o}}1 {Ã}{{\~A}}1 {Õ}{{\~O}}1
    {æ}{{\ae}}1 {Æ}{{\AE}}1 {ç}{{\c c}}1 {Ç}{{\c C}}1
    {ø}{{\o}}1 {Ø}{{\O}}1 {å}{{\r a}}1 {Å}{{\r A}}1
    {œ}{{\oe}}1 {Œ}{{\OE}}1 {ß}{{\ss}}1
    {ñ}{{\~n}}1 {Ñ}{{\~N}}1
}


% PROMPT 0
\textbf{Prompt 0 (Without class definitions)}
\begin{lstlisting}
Este modelo está diseñado para analizar tuits y clasificarlos en las siguientes categorías. Por favor, siga el formato de respuesta y utilice las etiquetas proporcionadas para cada categoría:

Categorías:
0 para Desacreditar Víctima
1 para Desacreditar Acto
2 para Insultos
3 para Desacreditar Deportista Autora

[Instrucciones detalladas sobre el análisis a realizar]
Evalúe el tuit en función del uso de palabras clave y frases asociadas a la incitación al odio, el contexto del mensaje, la intención que subyace en el mensaje, el impacto potencial del mensaje, el refuerzo de estereotipos o prejuicios de género, la presencia de amenazas o incitación a la violencia y el análisis del contenido multimedia. Tener en cuenta las políticas de incitación al odio de la plataforma y la legislación aplicable en materia de incitación al odio y discriminación. El objetivo es identificar mensajes que puedan ser ofensivos, discriminatorios o incitar al odio contra las mujeres en contextos deportivos, contribuyendo a un entorno seguro e inclusivo en el deporte.

Para cada tuit, facilite la clasificación de acuerdo con estas directrices y utilice un solo número. Ejemplos:

Tweet: "@Jennihermoso Vete a un Juzgado entonces. No lo harás porque sabes que es mentira." Response: {"prediction":"0"}

Tweet: "@Jennihermoso Jenni con todo el respeto del mundo, ¿agresión? Por favor..." Response: {"prediction":"1"}

Tweet: "@Jennihermoso Sinvergüenza!!! Mentirosa, jugar con algo tan serio. Se te tendría que caer la cara de vergüenza" Response: {"prediction":"2"}

Tweet: "@alexiaputellas @Jennihermoso Quieres otro balón de oro regalado eh." Response: {"prediction":"3"}
\end{lstlisting}
\newpage

% PROMPT 1
\textbf{Prompt 1 (With class definitions)}
\begin{lstlisting}
Este modelo está diseñado para analizar tuits y clasificarlos en las siguientes categorías. Por favor, siga el formato de respuesta y utilice las etiquetas proporcionadas para cada categoría:

Categorías:
0 para Desacreditar Víctima
1 para Desacreditar Acto
2 para Insultos
3 para Desacreditar Deportista Autora

Desacreditar Víctima: Comentarios que se centran en dañar la reputación de la víctima (en este caso Jennifer Hermoso).

Desacreditar Acto: Información que se basa en desacreditar lo sucedido, quitarle importancia al acto del beso no consentido.

Insultos: El usuario/a usa un lenguaje ofensivo, abusivo y vulgar dirigido sistemáticamente a una persona o grupos de personas.

Desacreditar Deportista Autora: Comentarios que se centran en dañar la reputación de la persona que escribe el tuit.

[Instrucciones detalladas sobre el análisis a realizar].
Evalúe el tuit en función del uso de palabras clave y frases asociadas a la incitación al odio, el contexto del mensaje, la intención que subyace en el mensaje, el impacto potencial del mensaje, el refuerzo de estereotipos o prejuicios de género, la presencia de amenazas o incitación a la violencia y el análisis del contenido multimedia. Tener en cuenta las políticas de incitación al odio de la plataforma y la legislación aplicable en materia de incitación al odio y discriminación. El objetivo es identificar mensajes que puedan ser ofensivos, discriminatorios o incitar al odio contra las mujeres en contextos deportivos, contribuyendo a un entorno seguro e inclusivo en el deporte.

Para cada tuit, facilite la clasificación de acuerdo con estas directrices y utilice un solo número. Ejemplos:

Tweet: "@Jennihermoso Vete a un Juzgado entonces. No lo harás porque sabes que es mentira." Response: {"prediction":"0"}

Tweet: "@Jennihermoso Jenni con todo el respeto del mundo, ¿agresión? Por favor..." Response: {"prediction":"1"}

Tweet: "@Jennihermoso Sinvergüenza!!! Mentirosa, jugar con algo tan serio. Se te tendría que caer la cara de vergüenza" Response: {"prediction":"2"}

Tweet: "@alexiaputellas @Jennihermoso Quieres otro balón de oro regalado eh." Response: {"prediction":"3"}
\end{lstlisting}
\newpage
